% \documentclass{beamer}
%Для защит онлайн лучше использовать разрешение 16x9
\documentclass[aspectratio=169]{beamer}

%%% PDF settings
\pdfvariable minorversion 7 % Set PDF version to 1.7.

%%% Fonts and language setup.
\usepackage{polyglossia}
% Setup fonts.
\usepackage{fontspec}
\setmainfont{CMU Serif}
\setsansfont{CMU Sans Serif}
\setmonofont{CMU Typewriter Text}

\usepackage{microtype} % Add fancy-schmancy font tricks

\usepackage{xcolor} % Add colors support.

%% Math
\usepackage{amsmath, amsfonts, amssymb, amsthm, mathtools} % Advanced math tools.
\usepackage{thmtools}
\usepackage{unicode-math} % Allow TTF and OTF fonts in math and allow direct typing unicode math characters.
\unimathsetup{
    warnings-off={
            mathtools-colon,
            mathtools-overbracket
        }
}
\setmathfont{Latin Modern Math} % default
\setmathfont[range={\setminus,\varnothing,\smashtimes}]{Asana Math}

%%% Images
\usepackage{graphicx}
\graphicspath{{figures/}}
\usepackage{import}

%%% Polyglossia setup after (nearly) everything as described in documentation.
\setdefaultlanguage{russian}
\setotherlanguage{english}

\usepackage{csquotes}

%%% Custom commands
\newcommand{\R}{\mathbb{R}}
\newcommand{\N}{\mathbb{N}}
\newcommand{\Z}{\mathbb{Z}}
\newcommand{\Q}{\mathbb{Q}}
\newcommand{\C}{\mathbb{C}}
\newcommand{\id}{\mathrm{id}}
\AtBeginDocument{\renewcommand{\leq}{\leqslant}}
\AtBeginDocument{\renewcommand{\geq}{\geqslant}}
\AtBeginDocument{\renewcommand{\Re}{\operatorname{Re}}}
\AtBeginDocument{\renewcommand{\Im}{\operatorname{Im}}}
\AtBeginDocument{\renewcommand{\phi}{\varphi}}
\AtBeginDocument{\renewcommand{\epsilon}{\varepsilon}}

%%% theorem-like envs
\theoremstyle{definition}

\declaretheoremstyle[spaceabove=0.5\topsep,
    spacebelow=0.5\topsep,
    headfont=\bfseries\sffamily,
    bodyfont=\normalfont,
    headpunct=.,
    postheadspace=5pt plus 1pt minus 1pt]{myStyle}
\declaretheoremstyle[spacebelow=\topsep,
    headfont=\bfseries\sffamily,
    bodyfont=\normalfont,
    headpunct=.,
    postheadspace=5pt plus 1pt minus 1pt,]{myStyleWithFrame}
\declaretheoremstyle[spacebelow=\topsep,
    headfont=\bfseries\sffamily,
    bodyfont=\normalfont,
    headpunct=.,
    postheadspace=5pt plus 1pt minus 1pt,
    qed=\blacksquare]{myProofStyleWithFrame}

\usepackage[breakable]{tcolorbox}
\tcbset{sharp corners=all, colback=white}
% \tcolorboxenvironment{theorem}{}
% \tcolorboxenvironment{theorem*}{}
% \tcolorboxenvironment{axiom}{}
% \tcolorboxenvironment{assertion}{}
% \tcolorboxenvironment{lemma}{}
% \tcolorboxenvironment{proposition}{}
% \tcolorboxenvironment{corollary}{}
% \tcolorboxenvironment{definition}{}
% \tcolorboxenvironment{proofReplace}{toprule=0mm,bottomrule=0mm,rightrule=0mm, colback=white, breakable }

\declaretheorem[name=Теорема, style=myStyleWithFrame]{theorem}
\declaretheorem[name=Теорема, numbered=no, style=myStyleWithFrame]{theorem*}
\declaretheorem[name=Аксиома, sibling=theorem, style=myStyleWithFrame]{axiom}
\declaretheorem[name=Преположение, sibling=theorem, style=myStyleWithFrame]{assertion}
\declaretheorem[name=Лемма, style=myStyleWithFrame]{lemma}
\declaretheorem[name=Предложение, sibling=theorem, style=myStyleWithFrame]{proposition}
\declaretheorem[name=Следствие, numberwithin=theorem, style=myStyleWithFrame]{corollary}

\declaretheorem[name=Определение, style=myStyleWithFrame]{definition}
\declaretheorem[name=Свойство, style=myStyle]{property}
\declaretheorem[name=Свойства, numbered=no, style=myStyle]{propertylist}

\declaretheorem[name=Пример, style=myStyle]{example}
\declaretheorem[name=Замечание, numbered=no, style=myStyle]{remark}

\declaretheorem[name=Доказательство, numbered=no, style=myProofStyleWithFrame]{proofReplace}
\renewenvironment{proof}[1][\proofname]{\begin{proofReplace}}{\end{proofReplace}}
% \declaretheorem[name=Доказательство, numbered=no, style=myProofStyleWithFrame]{longProof}

%%% Memoir settings
\chapterstyle{ger}
\setlength{\headheight}{2\baselineskip}

%%% HyperRef
\usepackage{hyperref}


% То, что в квадратных скобках, отображается внизу по центру каждого слайда.
\title[Инфраструктура тестирования]{Разработка инфраструктуры для автоматизации проверки задач по курсу \enquote{Теория формальных языков}}

% То, что в квадратных скобках, отображается в левом нижнем углу.
\institute[СПбГУ]{}

% То, что в квадратных скобках, отображается в левом нижнем углу.
\author[Николай Пономарев]{Николай Алексеевич Пономарев, группа 21.Б10-мм}

\begin{document}
{
\setbeamertemplate{footline}{}
% Лого университета или организации, отображается в шапке титульного листа
\begin{frame}
	\includegraphics[width=1.4cm]{pictures/SPbGU_Logo.png}
	\vspace{-35pt}
	\hspace{-10pt}
	\begin{center}
		\begin{tabular}{c}
			\scriptsize{Санкт-Петербургский государственный университет} \\
			\scriptsize{Кафедра системного программирования}
		\end{tabular}
		\titlepage
	\end{center}

	\btVFill

	{\scriptsize
		% У научного руководителя должна быть указана научная степень
		\textbf{Научный руководитель:} доцент кафедры системного программирования, к.ф.-м.н., Григорьев С. В. \\
	}
	\begin{center}
		\vspace{5pt}
		\scriptsize{Санкт-Петербург\\
			2024}
	\end{center}

\end{frame}
}

\begin{frame}
	\frametitle{Введение}
	\begin{itemize}
		\item На СПбГУ читается курс по формальным языкам
		\item Домашние задания требуют реализации нетривиальных алгоритмов
		\item Тесты обязательны, но не всегда покрывают граничные случаи
		\item Алгоритмические задачи хорошо проверяются с помощью property-based тестирования
	\end{itemize}
\end{frame}

% Обязательный слайд: четкая формулировка цели данной работы и постановка задачи
% Описание выносимых на защиту результатов, процесса или особенностей их достижения и т.д.
\begin{frame}
	\frametitle{Постановка задачи}
	\textbf{Цель}: разработать инфраструктуру для автоматизации проверки задач по курсу \enquote{Теория формальных языков}

	\textbf{Задачи}:
	\begin{enumerate}
		\item Изучить способы автоматического тестирования учебных задач
		\item Создать базовую инфраструктуру для автоматического тестирования
		\item Произвести проверку автоматических тестов, путём реализации решений
		\item Исправить накопленный технический долг
		      \begin{enumerate}
			      \item Сменить систему управления зависимостями
			      \item Обновить зависимости
		      \end{enumerate}
	\end{enumerate}
\end{frame}

\begin{frame}
	\frametitle{Виды автоматических тестов}

	Автоматические тесты можно характеризовать несколькими способами:
	\begin{itemize}
		\item По видимости
		      \begin{itemize}
			      \item открытые
			      \item закрытые
		      \end{itemize}
		\item По способу реализации
		      \begin{itemize}
			      \item ручные
			      \item property-based
			      \item смешанные
		      \end{itemize}
		\item По способу сдачи решений
		      \begin{itemize}
			      \item Learning Management Systems
			      \item Pull Request
			      \item \GitHub{} \textsc{Classroom}
		      \end{itemize}
	\end{itemize}
\end{frame}

\begin{frame}
	\frametitle{Требования к реализации}

	\begin{itemize}
		\item Тесты должны быть открытыми
		\item Тесты должны быть property-based
		\item Задачи должны сдаваться с помощью PR
		\item Имена и сигнатуры требуемых функций и классов известны заранее
		\item Архитектуру проекта и связь между модулями студенты определяют сами
	\end{itemize}

\end{frame}

\begin{frame}
	\frametitle{Используемые технологии}

	Используемые технологии во многом были зафиксированы
	\begin{itemize}
		\item \python{}~--- основной язык программирования
		\item \textsc{Pytest}~--- библиотека для тестирования
		\item \cfpqdata{}~--- библиотека-датасет графов и грамматик
		\item \pyformlang{}~--- библиотека для работы с грамматиками, регулярными выражениями и различными конечными автоматами
		\item \networkx{}~--- библиотека для работы с графами
		\item \scipy{}~--- библиотека для работы с разреженными матрицами
		\item \antlr{}~--- инструмент для генерации парсеров
		\item \textsc{Grammarinator}~--- инструмент генерации слов по грамматике \antlr{}
	\end{itemize}

\end{frame}

\begin{frame}
	\frametitle{Инфраструктура решения}

	Тесты к нереализованным студентом задачам необходимо пропускать.
	Это можно сделать с помощью:
	\begin{itemize}
		\item создания отдельных репозиториев для каждой задачи
		\item конфигурации системы тестирования
		\item проверки наличия требуемых функций
	\end{itemize}

\end{frame}

\begin{frame}[fragile=singleslide]
	\frametitle{Условный запуск тестов}

	\begin{center}
		\inputminted[]{python3}{pictures/minimal_example.py}
	\end{center}

\end{frame}

\begin{frame}[t]
	\frametitle{Структура решения}

	\begin{center}
		\includegraphics[height = 0.8\textheight, keepaspectratio, trim=0.8cm 0.8cm 0.8cm 0.8cm,clip]{pictures/project_structure.pdf}
	\end{center}
\end{frame}

\begin{frame}
	\frametitle{Обновление используемых инструментов}

	\begin{itemize}
		\item При реализации потребовалось обновить пакеты
		\item Получили проблемы с \pip{}
		\item В семестре обновили зависимости библиотеки \cfpqdata
		\item Перешли на \poetry
	\end{itemize}

\end{frame}

\begin{frame}
	\frametitle{Апробация}

	Весной курс читался первокурсникам магистратуры ИТМО.
	В течение семестра были выявлены следующие проблемы:
	\begin{itemize}
		\item Отсутствие ручных тестов
		\item Необходимость тестирования всех требуемых функции
		\item Необходимость предупреждать студентов о том, что решения нужно оптимизировать
	\end{itemize}

	\begin{figure}
		\caption{Зависимость задач друг от друга}
		\begin{center}
			\begin{tikzpicture}[shorten >=1pt, node distance=2cm, on grid]
				\node [circle, draw] (3) {3};
				\node [circle, draw, right = of 3] (4) {4};
				\node [circle, draw, right = of 4] (6) {6};
				\node [circle, draw, right = of 6] (7) {7};
				\node [circle, draw, right = of 7] (8) {8};
				\node [circle, draw, right = of 8] (9) {9};

				\path[->] (9) edge (8);
				\path[->] (8) edge (7);
				\path[->] (7) edge (6);
				\path[->] (6) edge (4);
				\path[->] (4) edge (3);

				\path[->] (7) edge [bend right] (4);
				\path[->] (8) edge [bend right] (4);
				\path[->] (9) edge [bend right] (4);

				\path[->] (8) edge [bend left] (6);
				\path[->] (9) edge [bend left] (6);

				\path[->] (9) edge [bend right] (7);
			\end{tikzpicture}
		\end{center}
	\end{figure}

\end{frame}

\begin{frame}
	\frametitle{Результаты}

	\begin{enumerate}
		\item Изучены способы автоматического тестирования учебных задач
		\item Создана базовая инфраструктура для автоматического тестирования
		\item Произведена проверка автоматических тестов, путём реализации решений
		\item Исправлен накопленный технический долг
	\end{enumerate}
	\blfootnote{Исходный код тестов и инфраструктуры: \url{https://github.com/FormalLanguageConstrainedPathQuerying/formal-lang-course}}
	\blfootnote{Имя пользователя: \texttt{WoWaster}}
	\blfootnote{Исходный код решений находится в приватном репозитории}
	\blfootnote{Pull Request в проект \cfpqdata{}: \url{https://github.com/FormalLanguageConstrainedPathQuerying/CFPQ_Data/pull/92}}

\end{frame}

\appendix

\begin{frame}
	\frametitle{Ход работы студента}

	\begin{center}
		\includegraphics[width=\textwidth, keepaspectratio]{pictures/student_workflow.pdf}
	\end{center}

\end{frame}

\begin{frame}
	\frametitle{Планы на будущее}

	Для исправления выявленных недостатков в будущем планируется:
	\begin{enumerate}
		\item добавить ручные тесты
		\item добавить тесты на наличие вспомогательных функций
		\item добавить ограничение на скорость исполнения алгоритмов
	\end{enumerate}

\end{frame}

\end{document}
